\input{preamble}

\section*{Annihilation 2}

From the first section we have
\begin{align*}
\sum_{ab}\mathcal M_1{}^{\mu\nu}\mathcal M^*_{1\mu\nu}
&=\frac{e^4}{(t-m^2)^2}
\left(8 t u - 24 m^2 t - 8 m^2 u - 8 m^4\right)
\\
\sum_{ab}\mathcal M_1{}^{\mu\nu}\mathcal M^*_{1\nu\mu}
&=\frac{e^4}{(t-m^2)(u-m^2)}
\left(8 m^2 s - 32 m^4\right)
\\
\sum_{ab}\mathcal M_2{}^{\nu\mu}\mathcal M^*_{2\nu\mu}
&=\frac{e^4}{(u-m^2)^2}
\left(8 t u - 8 m^2 t - 24 m^2 u - 8 m^4\right)
\end{align*}

We also have by the Casimir trick
\begin{align*}
\sum_{ab}\mathcal M_1{}^{\mu\nu}\mathcal M^*_{1\mu\nu}
&=\frac{e^4}{(t-m^2)^2}\operatorname{Tr}
\left(
%(\slashed p_1+m)\gamma^\mu(\slashed q_1+m)\gamma^\nu(\slashed p_2-m)\gamma_\nu(\slashed q_1+m)\gamma_\mu
(\slashed q_1+m)\gamma^\mu(\slashed p_2-m)\gamma_\mu(\slashed q_1+m)\gamma^\nu(\slashed p_1+m)\gamma_\nu
\right)
\\
\sum_{ab}\mathcal M_1{}^{\mu\nu}\mathcal M^*_{2\nu\mu}
&=\frac{e^4}{(t-m^2)(u-m^2)}\operatorname{Tr}
\left(
%(\slashed p_1+m)\gamma^\mu(\slashed q_2+m)\gamma^\nu(\slashed p_2-m)\gamma_\mu(\slashed q_1+m)\gamma_\nu
(\slashed q_2+m)\gamma^\mu(\slashed p_2-m)\gamma^\nu(\slashed q_1+m)\gamma_\mu(\slashed p_1+m)\gamma_\nu
\right)
\\
\sum_{ab}\mathcal M_2{}^{\nu\mu}\mathcal M^*_{2\nu\mu}
&=\frac{e^4}{(u-m^2)^2}\operatorname{Tr}
\left(
%(\slashed p_1+m)\gamma^\mu(\slashed q_2+m)\gamma^\nu(\slashed p_2-m)\gamma_\nu(\slashed q_2+m)\gamma_\mu
(\slashed q_2+m)\gamma^\mu(\slashed p_2-m)\gamma_\mu(\slashed q_2+m)\gamma^\nu(\slashed p_1+m)\gamma_\nu
\right)
\end{align*}

The trace operator Tr contracts (sums over) the first and last index.
Slash matrices are formed as
\begin{equation*}
\slashed p=p^\mu g_{\mu\nu}\gamma^\nu
=p^0\gamma^0-p^1\gamma^1-p^2\gamma^2-p^3\gamma^3
\end{equation*}

\href{https://georgeweigt.github.io/examples/annihilation-2-demo.html}{Eigenmath script}

\end{document}
